\ReportUnnumberedChapter{附录}
\label{chap:appendix}

\section*{技术资料与来源说明}
本报告的产品事实以当前项目源码、测试和归档输出为首要来源；需求分析及软件设计历史提交包仅用于追溯，若与当前实现不同，以当前源码及明确版本说明为准。NABCD用于组织需求、方法、收益、替代方案和交付渠道的定性分析，不据此宣称已经完成市场调研。以下官方资料用于解释术语与机制，项目实际采用范围仍以源码为准。
\begin{enumerate}
  \item Node.js 官方文档：HTTP、Crypto、SQLite 与测试框架。参考入口为 \url{https://nodejs.org/docs/latest-v24.x/api/http.html}、\url{https://nodejs.org/docs/latest-v24.x/api/crypto.html}、\url{https://nodejs.org/docs/latest-v24.x/api/sqlite.html}、\url{https://nodejs.org/docs/latest-v24.x/api/test.html}。项目使用版本见本轮环境记录，最新文档不代表已使用所有能力。
  \item SQLite 官方文档：事务与预写日志，\url{https://www.sqlite.org/lang_transaction.html}、\url{https://www.sqlite.org/wal.html}。报告中的数据结构来自\texttt{lib/db.mjs}，而非由官方示例推断。
  \item OMG Unified Modeling Language 2.5.1：\url{https://www.omg.org/spec/UML/2.5.1/}。本项目业务实现为函数式ES模块，用例图和实体关系图用于解释设计，不伪造面向对象实现类。
  \item 项目原始来源：\texttt{README.md}、\texttt{VERIFICATION.md}、\texttt{package.json}、\texttt{server.mjs}、\texttt{lib/}、\texttt{public/}及\texttt{tests/}。历史文档位于\texttt{需求分析/}，非功能测试工具位于\texttt{scripts/}。
  \item 场景与人物素材制作说明：项目\texttt{docs/watercolor-assets.md}记录既有素材及提示词。报告截图使用隔离账号和项目自带插画，不使用真实用户私有图片，也不把“项目自带”表述成公版或另行获得自由授权。
\end{enumerate}

\section*{附录A：证据目录与复现口径}
本附录以2026年11月4日当前源码为依据，列出可定位的实现与断言，并归档带日期的历史性能。节选是从真实源文件复制的连续片段，保留原变量、控制流和SQL，不是伪代码，不包含整份应用，也不应脱离导入、调用方和数据库结构独立运行。代码行号对应本次核查的源文件，不保证后续编辑后仍不变。节选中未加入隐去分支的省略号；外部依赖与节选范围在说明中标明。

\begin{table}[H]
  \centering
  \caption{证据索引与可复核内容}
  \label{tab:appendix-evidence}
  {\songti\rmfamily\zihao{5}
  \begin{tabular}{@{}p{.36\textwidth}p{.54\textwidth}@{}}
    \toprule
    文件或来源 & 核验内容 \\
    \midrule
    第02章需求与第04章矩阵 & FR-01至FR-15、IT-01至IT-15与实际测试文件的映射，编号为报告整理。 \\
    \texttt{package.json} & test仅运行\texttt{tests/*.test.mjs}；benchmark、展示与统计命令独立。 \\
    \texttt{tests/} & 规则、数据库与HTTP断言；另有浏览器脚本，不能由npm test的通过数代替。 \\
    \texttt{scripts/count-source.mjs} & 逐文件物理行/非空行、分类合计及排除清单。 \\
    \texttt{VERIFICATION.md} & 2026-10-31历史非功能结果、性能、运行环境及工作区状态；更早视觉记录单列日期。 \\
    \texttt{assets/report/regression.txt} & 2026-11-04，Windows/Node.js v24.13.0：120通过、0失败、3049.829 ms；取消、跳过和待办均为零。 \\
    \texttt{assets/report/source-count.txt} & 2026-11-04统计归档：后端3159、前端3061、CSS 4614、HTML 4、测试1960物理行；与第04章合计一致。 \\
    \texttt{assets/report/}五张PNG & 2026-11-04隔离场景截图生成及图像检查通过；不替代完整浏览器脚本的交互断言。 \\
    Git最近五条实际记录 & 第04章表\ref{tab:implementation-git-history}列出提交、作者日期、作者和原始标题；当前工作区并非干净状态。 \\
    \bottomrule
  \end{tabular}}
\end{table}

以下命令在项目根目录使用PowerShell执行。\texttt{npm test}开展隔离业务回归，本轮运行结果已有原始输出归档；统计结果另有逐文件归档。非功能展示另建运行目录和SQLite测试库，不读取日常\texttt{.env}，上游全部模拟，本轮未重跑，相关结果仍采用2026-10-31历史记录。浏览器需要预先配置测试工具及Chrome，具体命令与环境以各脚本和\texttt{VERIFICATION.md}为准；截图生成通过不等于这些独立浏览器脚本全部执行通过。

\begin{lstlisting}[language={},basicstyle=\ttfamily\small]
node --version
git rev-parse HEAD
git status --short
git log -5 --date=iso-strict --format="%h | %ad | %an | %s"
npm test
npm run count:source
npm run count:source -- --json
npm run showcase:nonfunctional
npm run showcase:nonfunctional -- --section=security
npm run showcase:nonfunctional -- --section=errors
npm run showcase:nonfunctional -- --section=robustness
npm run benchmark
\end{lstlisting}

历史完整展示结果记录于\texttt{runtime/nonfunctional/}下的运行目录\texttt{2026-10-31T08-20-15-184Z-25a5a154/}，\texttt{results.json}保存结构化数据，\texttt{summary.txt}为中文摘要，\texttt{app.db}是独立测试库。该路径来自验证记录，本次未据此假定所有忽略目录文件均可取得。缺失时应重新隔离运行并保留新日期结果，不能补造原日志。结果环境字段包括Node版本、平台、CPU、Git版本、工作区状态、数据库与模拟上游口径，能够防止把不同环境的测量混合。

\section*{附录B：关键实现的真实代码节选}
以下列表采用\texttt{listings}环境，\texttt{language=\{\}}表示不依赖额外的JavaScript语法配置。各节注明真实源文件、连续节选范围及其在完整调用链中的作用。

\subsection*{B.1 密码与会话摘要（FR-01）}
来源：\texttt{lib/auth.mjs}第35--53行。密码使用调用方传入的账号独立盐，scrypt产生64字节摘要；会话令牌为32随机字节的URL安全编码，仅存SHA-256摘要。完整认证还需要同文件的\texttt{verifyPassword()}、\texttt{login()}及\texttt{resolveSession()}，摘要本身不替代权限检查。
\begin{lstlisting}[language={},basicstyle=\ttfamily\small]
export function hashPassword(password, salt) {
  return scryptSync(password, salt, 64).toString('hex');
}

/**
 * 生成供浏览器持有的高熵会话令牌。
 * 使用 URL 安全编码，避免普通 Base64 中的特殊字符进入 Cookie 值。
 */
export function newToken() {
  return randomBytes(32).toString('base64url');
}

/**
 * 将会话令牌转换为固定长度 SHA-256 摘要。
 * 保存和查找会话均使用此摘要，原始令牌只交给浏览器。
 */
export function hashToken(token) {
  return createHash('sha256').update(String(token)).digest('hex');
}
\end{lstlisting}

\subsection*{B.2 双向匹配与事务（FR-02、FR-03）}
资料保存边界（源码核查）：\texttt{lib/matching.mjs}的\texttt{saveMatchingProfile()}只接收两段介绍，SQL写入介绍、标签状态与时间，不保存同意字段。注册/资料路由的请求同意检查不是持久授权记录；\texttt{analyzeProfile()}返回后比较介绍是否改变，也不是授权撤回或授权版本校验。IT-02中的资料并发变更和授权生命周期应列为待测/能力边界，不能从标签提取测试推断已实现完整同意管理。

来源：\texttt{lib/matching.mjs}第164--174行。两方向都必须有交集，平均值取整。调用前的固定词表、非空数量及去重检查由资料解析完成，不能向此函数传入任意未校验对象。
\begin{lstlisting}[language={},basicstyle=\ttfamily\small]
export function matchScore(a, b) {
  // 分别计算双方期望被满足的比例；任一方向没有交集时不予匹配。
  const towardB =
    a.desiredTags.filter((t) => b.selfTags.includes(t)).length /
    a.desiredTags.length;
  const towardA =
    b.desiredTags.filter((t) => a.selfTags.includes(t)).length /
    b.desiredTags.length;
  if (!towardA || !towardB) return 0;
  return Math.round((100 * (towardA + towardB)) / 2);
}
\end{lstlisting}

来源：\texttt{lib/draw.mjs}第224--234行。同步立即写事务仅在回调全部成功后提交，错误回滚。\texttt{draw()}、\texttt{decide()}和\texttt{remove()}共用此封装；不是在外部模型等待期间长期持有该写事务。
\begin{lstlisting}[language={},basicstyle=\ttfamily\small]
function transaction(db, fn) {
  db.exec('BEGIN IMMEDIATE');
  try {
    const value = fn();
    db.exec('COMMIT');
    return value;
  } catch (e) {
    db.exec('ROLLBACK');
    throw e;
  }
}
\end{lstlisting}

来源：\texttt{lib/draw.mjs}第272--281行，位于\texttt{draw()}内部，\texttt{pool}是已筛选候选，\texttt{pick}默认使用\texttt{randomInt(total)}，测试可注入票号。该连续片段只展示加权选择；后续同事务写入结果快照、前置待决定复用见完整源码。
\begin{lstlisting}[language={},basicstyle=\ttfamily\small]
const total = pool.reduce((n, p) => n + p.score, 0);
let ticket = pick(total),
  chosen = pool.at(-1);
for (const p of pool) {
  ticket -= p.score;
  if (ticket < 0) {
    chosen = p;
    break;
  }
}
\end{lstlisting}

\subsection*{B.3 真人消息参与者检查（FR-06）}
来源：\texttt{lib/real-chat.mjs}第242--261行。\texttt{conversation()}负责在写入前验证当前账号参与会话；\texttt{content()}限制正文。\texttt{sender\_id}由当前登录账号写入，不能由请求正文指定，更没有模型代替真人发言。纸条接受的事务建会话另见第188--236行。
\begin{lstlisting}[language={},basicstyle=\ttfamily\small]
export function sendRealMessage(
  db,
  userId,
  conversationId,
  body,
  now = Date.now()
) {
  const text = content(body, 500),
    row = conversation(db, userId, conversationId);
  const result = db
    .prepare(
      'INSERT INTO real_messages(conversation_id,sender_id,body,created_at) VALUES(?,?,?,?)'
    )
    .run(row.id, userId, text, now);
  db.prepare('UPDATE real_conversations SET updated_at=? WHERE id=?').run(
    now,
    row.id
  );
  return Number(result.lastInsertRowid);
}
\end{lstlisting}

\subsection*{B.4 情绪与资源回退（FR-09）}
来源：\texttt{public/portraits.mjs}第20--40行。\texttt{PORTRAIT\_EMOTIONS}为六种情绪白名单；置信度必须为70至100的整数。选择函数还核对资源存在；版本为零时由界面显示无图回退。情绪输入的消息编号校验由\texttt{generateGuidance()}完成，而非在该纯选择函数中完成。
\begin{lstlisting}[language={},basicstyle=\ttfamily\small]
export function portraitEmotion(analysis) {
  return PORTRAIT_EMOTIONS.some(([id]) => id === analysis?.label) &&
    Number.isInteger(analysis?.confidence) &&
    analysis.confidence >= 70 &&
    analysis.confidence <= 100
    ? analysis.label
    : 'default';
}

/**
 * 结合情绪判断与人物已上传的资源版本选择图片。
 * 对应情绪没有资源时回退默认版本，返回值供页面生成带版本号的地址。
 */
export function portraitSelection(person, analysis) {
  // 优先采用已上传的对应表情；缺少资源时回退到默认立绘版本。
  const emotion = portraitEmotion(analysis),
    version = person.portraitEmotions?.[emotion];
  return version
    ? { emotion, version }
    : { emotion: 'default', version: person.portraitVersion || 0 };
}
\end{lstlisting}

\subsection*{B.5 生图预览的版本校验与保存（FR-11）}
来源：\texttt{lib/portrait-generation.mjs}第120--134行，是\texttt{createPortraitGeneration()}的内部\texttt{apply()}。\texttt{find()}按账号与任务共同查找，\texttt{hash()}计算默认图摘要。就绪前不能保存；就绪后仍需默认图及情绪目标版本一致，保存图片与任务状态一起提交，随后清除预览字节。此实现并不省去用户对图片内容的人工确认。
\begin{lstlisting}[language={},basicstyle=\ttfamily\small]
function apply(userId,id){
 const job=find(userId,id);
 if(job.status==='saved')return {media:mediaStatus(db,userId),...state(userId)};
 if(job.status!=='ready')throw new ApiError('这张立绘尚不能保存。',409,'CONFLICT');

 // 默认图或目标情绪曾被替换时拒绝旧预览，避免覆盖用户更新的资源。
 const current=getMedia(db,userId,'portrait'),version=mediaStatus(db,userId).portraitEmotions[job.emotion]||0;
 if(!current||hash(current.bytes)!==job.source_hash||version!==job.target_version)throw new ApiError('默认立绘或该情绪图片已更换，请放弃旧预览后重新生成。',409,'IMAGE_STALE');
 db.exec('BEGIN IMMEDIATE');
 try{
  saveMedia(db,userId,'portrait','data:image/png;base64,'+Buffer.from(job.bytes).toString('base64'),job.emotion);
  db.prepare("UPDATE portrait_generation_jobs SET status='saved',bytes=NULL WHERE id=?").run(id);db.exec('COMMIT');
 }catch(error){db.exec('ROLLBACK');throw error}
 return {media:mediaStatus(db,userId),...state(userId)};
}
\end{lstlisting}

\section*{附录C：真实测试断言与检查定位}
\subsection*{C.1 待决定复用与同日继续抽取}
来源：\texttt{tests/draw.test.mjs}第26--35行。\texttt{fixture()}准备内存库和就绪账号，\texttt{now}为固定时间；断言不仅验证重复抽取，还覆盖跨日复用、重复喜欢、虚拟日期推进/重置不改变抽取口径，以及决定后同日新增记录。
\begin{lstlisting}[language={},basicstyle=\ttfamily\small]
test('unlimited same-day draws preserve pending results and idempotent decisions',t=>{
 const {db,a}=fixture(t),one=draw(db,a.user.id,now,()=>0);assert.equal(one.remaining,null);assert.equal(one.unlimited,true);
 assert.equal(draw(db,a.user.id,now,()=>100).pending.drawId,one.pending.drawId);
 assert.equal(draw(db,a.user.id,now+86400000).pending.drawId,one.pending.drawId);
 decide(db,a.user.id,one.pending.drawId,'like',now);decide(db,a.user.id,one.pending.drawId,'like',now);
 assert.equal(drawState(db,a.user.id,now).collection.length,1);assert(drawState(db,a.user.id,now).canDraw);
 runAction(db,a.user,'advance',{days:30},now);runAction(db,a.user,'reset',{},now);
 assert.equal(drawState(db,a.user.id,now).remaining,null);assert.equal(drawState(db,a.user.id,now).collection.length,1);
 const two=draw(db,a.user.id,now,()=>0);assert.notEqual(two.pending.drawId,one.pending.drawId);assert.equal(db.prepare('SELECT COUNT(*) AS n FROM draw_records WHERE user_id=? AND day=?').get(a.user.id,drawDay(now)).n,2);
});
\end{lstlisting}

\subsection*{C.2 私有预览、越权拒绝与幂等保存}
来源：\texttt{tests/portrait-generation.test.mjs}第24--30行，位于“generation uses stored own reference, remains a private preview, and saves only on confirmation”测试内部；\texttt{fixture()}在此前注入模拟上游PNG。\texttt{a}为所有者，\texttt{b}为其他账号，\texttt{finished()}轮询任务完成。连续断言展示请求幂等、预览阶段未写目标槽、越权404、透明像素、确认后saved及字节释放。
\begin{lstlisting}[language={},basicstyle=\ttfamily\small]
const request=input(),job=service.start(a,request);assert.equal(service.start(a,request).id,job.id);assert.equal(calls,1);
assert.throws(()=>service.start(a,input()),e=>e.code==='IMAGE_PENDING');
assert.equal((await finished(service,a)).status,'ready');assert.equal(mediaStatus(db,a).portraitEmotions.happy,undefined);
assert.throws(()=>service.image(b,job.id),e=>e.status===404);assert.throws(()=>service.apply(b,job.id),e=>e.status===404);
assert(decodePortraitPng(Buffer.from(service.image(a,job.id))).pixels[3]===0);
assert(service.apply(a,job.id).media.portraitEmotions.happy);assert.equal(service.apply(a,job.id).job.status,'saved');
assert.equal(db.prepare('SELECT bytes FROM portrait_generation_jobs WHERE id=?').get(job.id).bytes,null);
\end{lstlisting}

\subsection*{C.3 其他关键测试的定位}
认证的\texttt{logging in on a second device kicks the first one off}测试同时检查旧会话\texttt{replaced}、新会话有效及库内有效会话数为一；\texttt{core.test.mjs}中的\texttt{replies retry without duplicate user message, both directions}检查重复发送拒绝、回复只增一条及重复完成返回false。\texttt{mutual outcome requires current explicit consent and spends once}与撤回/屏蔽用例检查跨期重新授权、不合资格不耗次数。

\texttt{server.test.mjs}的\texttt{private data is excluded from upstream messages}核对上游不含私人分数、授权字段及密钥；\texttt{full AI capacity still permits profile saves and manual messages}使用挂起模型请求验证普通业务与文本四并发额度分离。\texttt{favorability.test.mjs}中的\texttt{privacy switch: reciprocal visibility enforcement}及HTTP API用例核对双方开关与参与者权限。媒体分离测试验证旧头像迁移只执行一次、保留历史资源且不挪用作新立绘。这些定位与第04章矩阵共同使用，不把测试标题本身当作执行日志。

\subsection*{C.4 断言范围与源码边界}
第04章每个IT场景均分别列出已有断言、源码与待测范围；同一例包含多个层次，不设整体“通过”状态。已有断言必须能在对应测试中定位，120项是归档框架统计，不能证明所有IT子场景已执行。历史非功能展示的五并发和失败重试不能当作本轮新增断言；当前抽取HTTP测试实际发出八个并发请求。

\texttt{actions.test.mjs}的reset测试只断言虚拟纸条清空、偏移归零、展示名回到用户名及账号存续；\texttt{draw.test.mjs}另断言reset后收藏数量不变。源码\texttt{ACTIONS.reset}调用\texttt{freshState()}、\texttt{defaultProfile()}并覆盖虚拟状态，\texttt{writeState()}只更新users资料与states；不删除真人纸条/会话/消息、收藏/抽取、匹配资料或媒体。真人和媒体完整保留的专门断言、页面二次确认及在途请求与重置竞态仍属待测，不把测试标题中的“interactions”扩张为全部业务数据。

真人参考分的隐私测试采用顺序切换：先开启并读取/分析，完成后关闭再读取，核对隐藏；没有在模型挂起期间关闭任一方开关的断言。源码\texttt{/api/real/favorability/analyze}在模型请求前读取双方开关，但\texttt{await completion()}之后不复查，直接写\texttt{conversation\_favorability}并返回\texttt{canView:true}。在途关闭不能撤回已外发内容，且原响应仍可能给出结果；后续GET隐藏不等于在途撤回保护。该缺口属于源码发现，动态竞态验证待测，未因本轮120项通过而消失。

\section*{附录D：历史性能明细（2026-10-31）}
以下全部数字抄录\texttt{VERIFICATION.md}的2026年10月31日记录，环境为Core Ultra 7 155H、WSL2 Linux 6.6.87.2、Node.js v24.21.0，数据库位于E盘独立测试目录，SQLite文件库启用WAL。每组预热20次、正式100次，准备及预热不计时；每组正式执行前重置初始消息规模，发送组再新增100条。客户端与服务端为不同进程，使用模拟模型，未测真实推理、生图或公网延迟。表中“读”指真人快照，不是无限历史分页；每行错误均为0/100。

\begin{table}[H]
  \centering
  \caption{2026-10-31 历史性能：并发1}
  \label{tab:appendix-performance-one}
  {\songti\rmfamily\zihao{5}
  \begin{tabular}{@{}rlrrrrr@{}}
    \toprule
    初始消息 & 接口 & 平均/ms & P50/ms & P95/ms & 成功请求/秒 & 错误 \\
    \midrule
    100 & 状态 & 0.83 & 0.72 & 1.37 & 1206.20 & 0/100 \\
    100 & 账号 & 4.33 & 4.31 & 5.30 & 230.83 & 0/100 \\
    100 & 读 & 5.78 & 5.77 & 6.88 & 172.94 & 0/100 \\
    100 & 发 & 13.25 & 13.09 & 15.24 & 75.45 & 0/100 \\
    1000 & 状态 & 0.40 & 0.38 & 0.57 & 2521.92 & 0/100 \\
    1000 & 账号 & 4.54 & 4.38 & 6.36 & 220.06 & 0/100 \\
    1000 & 读 & 5.65 & 5.52 & 6.96 & 176.95 & 0/100 \\
    1000 & 发 & 12.76 & 12.18 & 15.37 & 78.37 & 0/100 \\
    \bottomrule
  \end{tabular}}
\end{table}

\begin{table}[H]
  \centering
  \caption{2026-10-31 历史性能：并发10}
  \label{tab:appendix-performance-ten}
  {\songti\rmfamily\zihao{5}
  \begin{tabular}{@{}rlrrrrr@{}}
    \toprule
    初始消息 & 接口 & 平均/ms & P50/ms & P95/ms & 成功请求/秒 & 错误 \\
    \midrule
    100 & 状态 & 2.54 & 2.42 & 3.29 & 3725.35 & 0/100 \\
    100 & 账号 & 35.42 & 30.06 & 64.25 & 268.69 & 0/100 \\
    100 & 读 & 50.31 & 41.61 & 94.18 & 190.26 & 0/100 \\
    100 & 发 & 122.96 & 103.06 & 221.54 & 77.57 & 0/100 \\
    1000 & 状态 & 2.33 & 2.09 & 4.34 & 4160.91 & 0/100 \\
    1000 & 账号 & 31.45 & 27.21 & 57.56 & 302.86 & 0/100 \\
    1000 & 读 & 43.39 & 36.53 & 77.90 & 219.30 & 0/100 \\
    1000 & 发 & 111.58 & 92.20 & 201.49 & 85.55 & 0/100 \\
    \bottomrule
  \end{tabular}}
\end{table}

“状态、账号、读、发”分别对应\texttt{GET /api/status}、\texttt{GET /api/me}、\texttt{GET /api/real}与\texttt{POST /api/real/message}。16组正式请求合计1600次、无失败；平均与分位数只按成功请求计算，P50/P95按排序后最近秩取值，吞吐用成功数除以该批次墙钟时间，不用\texttt{1000/平均耗时}替代并发吞吐。历史非功能检查19项通过与历史回归120项通过属于另外两种统计；2026-11-04归档回归的3049.829 ms也不是此表中任一列。

\section*{附录E：本轮截图与待验证范围}
五张截图于2026年11月4日通过隔离场景生成，经PNG文件头核对，原图均为1440$\times$1000像素，存放于\texttt{assets/report/}。图像检查确认文件可读且呈现对应页面；表\ref{tab:appendix-screenshots}记录实际可见内容而不推断图外操作。截图采用演示场景，消息、情绪和参考值不构成真实心理或真实模型质量证据。

\begin{table}[H]
  \centering
  \caption{2026-11-04 五张隔离截图的检查索引}
  \label{tab:appendix-screenshots}
  {\songti\rmfamily\zihao{5}
  \begin{tabular}{@{}p{.15\textwidth}p{.40\textwidth}p{.35\textwidth}@{}}
    \toprule
    文件 & 已检查的可见内容 & 不能由该图单独确认的行为 \\
    \midrule
    \texttt{login.png} & 登录表单、账号单会话说明、注册入口。 & 注册成功、非法凭据拒绝、旧设备会话失效。 \\
    \texttt{meet.png} & 不限次数标识、暂无合适候选提示、档位说明及4/5收藏占位。 & 加权抽取、并发复用、喜欢/跳过/冷却与刷新恢复。 \\
    \texttt{notes.png} & 真人纸条、已回信状态、真人对话入口与虚拟人物分区。 & 角色权限、拒绝不建会话、接受的事务原子性及实时更新。 \\
    \texttt{chat.png} & 演示对象、消息、三项建议和自己输入、风格选项、开心标识及参考分。 & 建议仅填草稿、消息/风格竞态、轮询、失败恢复及模型真实性。 \\
    \texttt{media.png} & 媒体弹窗、背景/头像/默认立绘的独立入口、上传限制与当前预览。 & 上传持久化、全部情绪槽、生图授权、私有预览及确认保存。 \\
    \bottomrule
  \end{tabular}}
\end{table}

\subsection*{E.1 回归归档与读取范围}
\texttt{assets/report/regression.txt}包含\texttt{npm test}的测试标题、单项耗时和框架汇总；末尾汇总为tests 120、suites 0、pass 120、fail 0、cancelled 0、skipped 0、todo 0、duration\_ms 3049.829。Windows和Node.js v24.13.0为本轮记录的执行环境，不把单项测试耗时或总耗时作为接口性能指标。\texttt{source-count.txt}为统计命令的逐文件JSON输出（含npm命令头），其中主体业务6220行、应用源码10838行；其物理行和非空行均可能包含注释。

环境与工作区旁证保存在\texttt{assets/report/provenance.txt}，采集时间为2026-11-04 19:42（UTC+8），包含操作系统内核标识、PowerShell版本、Node.js版本、Git HEAD、非干净工作区列表，以及主体源码和回归文件的SHA-256摘要。它在回归之后采集，不能伪称回归日志本身内嵌了时间和版本；本次没有修改业务源文件，摘要用于后续判断文件是否已变化。HEAD为\texttt{9d18f30668073a988bf5e1181b8c288e370faddf}，不代表未提交内容已经进入该提交。

\subsection*{E.2 尚需独立验证的能力}
除浏览器复测外，IT各例仍有独立覆盖缺口：匹配资料在途变化与同意生命周期、真人纸条额度和拒绝/事务故障、真人模型故障下手动发送、重置后真人及媒体完整保留，以及参考分分析在途关闭。对已由源码确认不存在的持久同意或返回后开关复查，不能写成“已通过、仅待截图”；待测表示证据不足，不表示业务缺口已修复。

本轮已归档自动化业务回归、源码规模和五张页面截图，但完整浏览器回归尚未整套重跑。消息/风格竞态、建议只填草稿、会话失效后的草稿恢复、重置二次确认、生图授权与预览保存，以及移动端多视口、键盘焦点和减少动态效果仍需对应版本的脚本输出。真实模型兼容性、文本质量、人物一致性及授权外发载荷应另记服务配置范围与人工检查；更高负载、无障碍和生产安全不能由历史性能或截图推导。

历史2026-09-26的“没有真实多人通信”只描述当时版本，当前源码已实现账号间文字通信；当前仍没有真实身份核验、生产审核、端到端加密或公网生产保障。未执行的检查不填通过，提交历史不等于部署记录，行数和截图数量也不等于个人贡献量。
